% =====================================================================
%  Höhe des Langlochs und Anlegen immer von derselben Seite.
%  Schnitt durch zwei stumpf gestoßene Teile. Selbst gezeichnet.
% =====================================================================
\begin{tikzpicture}[x=1mm, y=1mm, font=\scriptsize]

  % Dübel zwischen zwei Teilen, Mitte bei y=#1 (Schnitt)
  \newcommand{\duebelschnitt}[2]{%
    \fill[nut] (#2+13,#1-2.5) rectangle (#2+26,#1+2.5) (#2+28,#1-2.5) rectangle (#2+41,#1+2.5);
    \draw[fill=holz!70!white, draw=holzrand, line width=0.5pt, rounded corners=1.2]
      (#2+14.5,#1-2.2) rectangle (#2+39.5,#1+2.2);}

  % ---------------- 1 Höhe -----------------------------------------
  \feld{0}{0}{54}{50}{1\quad Höhe einstellen}
  \draw[werkstueck] (6,12) rectangle (44,31);
  \fill[black!70, rounded corners=2.4] (16,19) rectangle (34,24);
  \draw[metall, fill=black!30] (4,31) rectangle (46,34);
  \node[anchor=south] at (25,34) {Fräsanschlag liegt oben auf};
  \draw[hilfslinie] (34,21.5) -- (50,21.5) (44,31) -- (50,31);
  \draw[mass] (48,21.5) -- (48,31);
  \node[anchor=west, fill=feld, inner sep=0.4pt] at (48.6,26.25) {h};
  \node[align=center] at (27,5.5) {meist halbe Materialstärke};

  % ---------------- 2 richtig --------------------------------------
  \feld{58}{0}{54}{50}{2\quad Gleiche Seite oben}
  \draw[werkstueck] (62,14) rectangle (84,28) (86,14) rectangle (108,28);
  \duebelschnitt{22.5}{58}
  \node[anchor=south] at (85,28.5) {Sichtseite};
  \draw[hilfslinie] (62,28) -- (108,28);
  \pic at (103,44) {ok};
  \node[align=center] at (85,5.5) {beide Teile von der Sichtseite\\angelegt: bündig};

  % ---------------- 3 falsch ---------------------------------------
  \feld{116}{0}{54}{50}{3\quad Seite gewechselt}
  \draw[werkstueck] (120,14) rectangle (142,28) (144,17) rectangle (166,31);
  \duebelschnitt{22.5}{116}
  \draw[hilfslinie] (138,28) -- (150,28);
  \draw[mass] (147,28) -- (147,31);
  \node[anchor=west] at (147.5,33) {Versatz};
  \pic at (161,44) {nok};
  \node[align=center] at (143,5.5) {Langloch nicht genau mittig und\\einmal von unten angelegt};
\end{tikzpicture}
